%% P26 \dash{} Wnioskowanie statystyczne dla inżyniera
%%
%% Zalążek. Poniższy opis jest kontraktem tego programu: co ma uzasadnić i co
%% ma dać czytelnikowi. Napisanie programu oznacza usunięcie bloku
%% \programstub{} i zastąpienie go programem. Jeśli po przerobieniu materiału
%% opis okaże się błędny, popraw go i odnotuj sprzeczność w CLAUDE.md w sekcji
%% *Resolved questions*.

\program{Wnioskowanie statystyczne dla inżyniera}
\label{prog:P26}

\programstub{%
\textit{[Opis do przetłumaczenia \dash{} patrz wydanie angielskie.]}

Argument: a measured difference is a random quantity, and the question is
whether it is larger than the noise in the measurement. Payoff: the program
that carries this house's first rule into a mathematics book. \enquote{Model B scored 71.4 and model A scored 70.9 on 200 evaluation items} \dash{} is that
real? Worked end to end: the standard error of a proportion; a bootstrap
confidence interval on an evaluation set; the paired comparison, because A
and B saw the same prompts and a paired test has far more statistical power;
what a p-value does and does not say, stated flatly; the multiple-
comparisons problem, which is what a leaderboard with forty models is; and a
power calculation answering how many evaluation items are needed to detect a
one-point difference at all. The honest conclusion \dash{} that most published
leaderboard deltas of this size are not distinguishable from noise, and the
reader can now check.

\medskip
\textbf{Planowana długość:} 60 ramek.
}
